\chapter{Conclusion}\label{ch:conclusion}

\paragraph{}As this project reaches its completion, Hubly can be considered to have fulfilled the main objectives defined at the beginning of the work. The project involved the development of a full-stack web application using .NET for the backend, PostgreSQL for the relational database, and React/Next.js for the web frontend. This technology stack presented several technical challenges, but it also provided an opportunity to understand the interaction between backend services, database design and frontend application development in a web application context.

Developing Hubly allowed us to apply and expand the theoretical knowledge acquired throughout the academic program, connecting computer engineering concepts with the design of a marketplace-oriented application. The project also provided experience with different development paradigms, such as the typed backend architecture of .NET and the component-based structure of React. Working with these technologies contributed to the development of full-stack skills, problem-solving abilities and a better understanding of how different system components interact within a complete application.

The implementation of Hubly required addressing several software engineering and architectural challenges. One relevant aspect was the design of the relational database schema, where a single user can act as a company or as a creator with multiple social profiles. Each social profile required its own metadata, messaging context and private conversation tags. Another important component was the implementation of multi-criteria search mechanisms for both \emph{CreatorSearchInputModel} and \emph{CompanySearchInputModel}, allowing users to filter results according to sectors, price ranges, audience size, company size and location. In addition, the internal messaging system was designed to associate conversations with specific company or creator social profiles, rather than only with generic user accounts. The system also includes profile-view history and recommendation logic, which provide a basis for suggesting companies or creator social profiles according to previous interactions.

Looking forward, Hubly could be extended in several ways. Future work may include the integration of social media APIs, such as Meta and YouTube APIs, to support audience verification, metric collection and demographic data retrieval. Another possible improvement would be the implementation of a multi-user collaboration system, allowing marketing or management teams to administer the same company or creator profile. To support this type of shared access, concurrency control mechanisms, such as optimistic locking, could be introduced to prevent conflicts when multiple users modify the same information. An audit logging system could also be added to record relevant user actions, including the user responsible for the action, the affected area of the application and the corresponding timestamp.

Overall, Hubly provides a structured solution for supporting interactions between companies and digital creators. By centralizing profile discovery, profile-based conversations and private negotiation tags, the system reduces dependence on scattered external communication channels. The final prototype demonstrates how a web application can support company-creator partnerships through organized search, communication and interaction tracking workflows.
